\setcounter{section}{4} % Setting counter for Week 5

\section{Optimization}

Let's look at finding the extreme values of a function $f: U \subseteq \R^n \to \R$. 

\begin{definition*}[Local Extrema and Critical Points]
\begin{itemize}
    \item $f$ has a \textbf{local minimum} at $x_0 \in U$ if there exists a ball $B_\eps(x_0)$ such that $\forall x \in B_\eps(x_0) \cap U$, we have $f(x) \ge f(x_0)$. \textcolor{eGray}{(Analogous for local maximum with $\le$).}
    \item Extrema are \textbf{strict} if the inequality is strict ($<$ or $>$).
    \item A \textbf{critical point} of $f \in C^1(U, \R)$ is any point $x_0$ where the gradient vanishes: $\nabla f(x_0) = 0$.
\end{itemize}
\end{definition*}

\begin{proposition*}
If $f \in C^1(U, \R)$ has a local extremum at an \emph{interior} point $x_0 \in \operatorname{int}(U)$, then $x_0$ must be a critical point ($\nabla f(x_0) = 0$).
\end{proposition*}

\gemininote{Geometrically, the gradient $\nabla f(x_0)$ always points in the direction of the \qt{steepest ascent}. If you are at a local minimum (the bottom of a valley) or a local maximum (the peak of a mountain), there is no \qt{uphill} direction anymore. You are standing perfectly flat. Therefore, the gradient vector must be the zero vector!}

\begin{center}
\begin{tikzpicture}[scale=1.5, very thick]
    % Coordinate grid
    \draw[eGray, dashed] (-1,-1) grid (3,2);
    \draw[->, eBlack] (-1.5, 0) -- (3.5, 0) node[right] {$x$};
    \draw[->, eBlack] (0, -1.5) -- (0, 2.5) node[above] {$y$};
    
    % Contour lines (Level sets)
    \draw[eMediumOrchid] (1,0.5) ellipse (1.5cm and 0.8cm);
    \draw[eMediumBlue] (1,0.5) ellipse (0.8cm and 0.4cm);
    
    % Minimum point
    \fill[eDarkGreen] (1,0.5) circle (2pt) node[below right] {Minimum};
    
    % Gradient vector
    \coordinate (P) at (2.2, 0.74); % Point on the outer contour
    \fill[eSaddleBrown] (P) circle (1.5pt);
    \draw[->, eDarkRed, line width=1.5pt] (P) -- ++(0.8, 0.6) node[above] {$\nabla f$};
    \node[eDarkRed, align=center] at (3.5, 1.6) {\small Steepest \\ \small ascent};
\end{tikzpicture}
\end{center}

\hrulefill

\section{Constrained Optimization (Lagrange)}

Suppose we want to find the hottest place on the surface of the Earth. 
In theory, temperature is a scalar field defined everywhere in space: $T \in C^\infty(\R^3, (0, \infty))$. We model the Earth's surface as the unit sphere $S^2 := \set{x \in \R^3 : \abs{x}^2 = 1}$.

We want to find the local maxima of the \emph{restricted} function $T|_{S^2}$. 
Crucially, these are \textbf{not necessarily} critical points of the global function $T$! 

\gemininote{Imagine standing on the equator. The true gradient of the temperature $\nabla T$ might point straight up into the sun (off the surface of the Earth). However, since you are constrained to walk only \emph{along} the surface, you only care about the \qt{shadow} of that gradient on the ground. A point is a constrained extremum if the true gradient points entirely off the surface, meaning it is perfectly perpendicular to the ground!}

Since our constraint surface $S^2$ is the level set of a function $g(x) = \abs{x}^2 - 1 = 0$, we know that $\nabla g(x)$ is everywhere perpendicular to the surface. 
Thus, at a constrained extremum $p$, the temperature gradient $\nabla T(p)$ must be perfectly parallel to the constraint gradient $\nabla g(p)$.

\begin{proposition*}[Lagrange Multipliers]
Suppose $f: U \subseteq \R^n \to \R$ and $g = (g_1, \dots, g_k): U \to \R^k$ are $C^1$. 
If the restricted function $f|_{g^{-1}\set{0}}$ has a local extremum at $x_0$, then the gradients $\set{\nabla f(x_0), \nabla g_1(x_0), \dots, \nabla g_k(x_0)}$ are linearly dependent.

Assuming the constraint gradients are independent, this means $\nabla f(x_0)$ is a linear combination of the constraint gradients. We can find these points by looking for critical points of the \textbf{Lagrangian}:
\[ L(x, \lambda) := f(x) - \sum_{i=1}^k \lambda_i g_i(x) \]
We force $\nabla L(x, \lambda) = 0$, which yields the system:
\begin{align*}
    \nabla f(x) - \lambda \cdot \nabla g(x) &= 0 \quad \text{\textcolor{eGray}{(Parallel gradients)}} \\
    g_j(x) &= 0 \quad \text{\textcolor{eGray}{(You are actually on the surface)}}
\end{align*}
\end{proposition*}

\begin{exercise*}
Find the local extrema of $f(x,y,z) = xyz$ on the sphere $K = \set{(x,y,z) \in \R^3 : x^2 + y^2 + z^2 \le 1}$.
\end{exercise*}

\begin{solution}
First, we check the \textbf{interior} ($x^2 + y^2 + z^2 < 1$) for standard critical points:
\[ \nabla f(x,y,z) = (yz, xz, xy) = 0 \implies \text{at least two variables must be } 0. \]
For example, $(0,0,0)$, or $(x,0,0)$. All these yield $f = 0$.

Next, we check the \textbf{boundary} using Lagrange multipliers. The constraint is $g(x,y,z) = x^2 + y^2 + z^2 - 1 = 0$. The Lagrangian is:
\[ L(x,y,z,\lambda) = xyz - \lambda(x^2 + y^2 + z^2 - 1) \]
Taking the partial derivatives and setting them to zero gives our system:
\begin{enumerate}
    \item $\partial_x L = yz - 2\lambda x = 0 \implies yz = 2\lambda x$
    \item $\partial_y L = xz - 2\lambda y = 0 \implies xz = 2\lambda y$
    \item $\partial_z L = xy - 2\lambda z = 0 \implies xy = 2\lambda z$
    \item $\partial_\lambda L = 1 - (x^2 + y^2 + z^2) = 0 \implies x^2 + y^2 + z^2 = 1$
\end{enumerate}

\gemininote{We want to divide by $x$, $y$, and $z$ to isolate $2\lambda$. Is that allowed? If $x=0$, equation (1) forces $y=0$ or $z=0$. If $x=0$ and $y=0$, equation (4) forces $z = \pm 1$. The points $(0,0,\pm 1)$ yield $f=0$, which we already know are not strict extrema since we can easily find positive/negative values elsewhere. Thus, we can safely assume $x, y, z \ne 0$ for the true extrema and divide!}

Assuming non-zero coordinates, we isolate $2\lambda$:
\[ 2\lambda = \frac{yz}{x} = \frac{xz}{y} = \frac{xy}{z} \]
Multiplying the terms pairwise yields:
\[ y^2 z^2 = x^2 z^2 = x^2 y^2 \implies x^2 = y^2 = z^2 \]
Plugging this symmetry into the constraint (4):
\[ 3x^2 = 1 \implies x = \pm \frac{1}{\sqrt{3}} \]
This gives $2^3 = 8$ solutions on the sphere:
\[ (x,y,z) = \left( \frac{\alpha}{\sqrt{3}}, \frac{\beta}{\sqrt{3}}, \frac{\gamma}{\sqrt{3}} \right) \quad \text{where } \alpha, \beta, \gamma \in \set{-1, 1} \]
\end{solution}

\hrulefill

\section{The Hessian Test}

Just like the second derivative test in Analysis 1, we can classify critical points in $\R^n$ using the second derivatives. The \textbf{Hessian matrix} of $f \in C^2(U, \R)$ is:
\[ \mathscr{H}f(x) = (\partial_i \partial_j f)_{i,j=1,\dots,n} \]
By Schwarz's Lemma, this matrix is perfectly symmetric.

\begin{definition*}[Definiteness of a Symmetric Matrix]
Let $A \in \R^{n \times n}$ be symmetric. It is called:
\begin{itemize}
    \item \textbf{Positive definite ($A > 0$):} $v^T A v > 0 \quad \forall v \ne 0$
    \item \textbf{Negative definite ($A < 0$):} $v^T A v < 0 \quad \forall v \ne 0$
    \item \textbf{Indefinite:} $\exists v, u$ such that $v^T A v > 0$ and $u^T A u < 0$
\end{itemize}
\end{definition*}

If $x_0$ is a critical point ($\nabla f(x_0) = 0$), then:
\begin{itemize}
    \item $\mathscr{H}f(x_0)$ is positive definite $\implies$ \textcolor{eDarkGreen}{\textbf{Local Minimum}}
    \item $\mathscr{H}f(x_0)$ is negative definite $\implies$ \textcolor{eDarkRed}{\textbf{Local Maximum}}
    \item $\mathscr{H}f(x_0)$ is indefinite $\implies$ \textcolor{eSaddleBrown}{\textbf{Saddle Point}}
\end{itemize}

\subsection*{How to determine Definiteness?}
Let $\lambda_1, \dots, \lambda_n$ be the eigenvalues of $A$. Recall that $\det(A) = \prod \lambda_i$ and $\operatorname{Tr}(A) = \sum \lambda_i$.
\begin{itemize}
    \item \textbf{For $n=2$:} 
        \begin{itemize}
            \item $\det(A) < 0 \implies$ Indefinite
            \item $\det(A) > 0$ and $\operatorname{Tr}(A) > 0 \implies$ Positive definite
            \item $\det(A) > 0$ and $\operatorname{Tr}(A) < 0 \implies$ Negative definite
        \end{itemize}
    \item \textbf{For $n \ge 3$ (The Hurwitz Criterion):} Check the leading principal minors (the determinants of the upper-left submatrices $A_1, A_2, \dots, A_n$).
        \begin{itemize}
            \item $A_1 > 0, A_2 > 0, A_3 > 0 \implies$ Positive definite
            \item $A_1 < 0, A_2 > 0, A_3 < 0 \implies$ Negative definite (alternating signs!)
            \item Anything else (excluding zeros) $\implies$ Indefinite
        \end{itemize}
\end{itemize}

\begin{example*}[Testing Matrices]
\textbf{1.} $A = \begin{pmatrix} 2 & 2 \\ 2 & 9 \end{pmatrix}$. $\det = 18 - 4 = 14 > 0$. $\operatorname{Tr} = 11 > 0 \implies$ \textcolor{eDarkGreen}{Positive definite}.

\textbf{2.} $A = \begin{pmatrix} 4 & 2 \\ 2 & -4 \end{pmatrix}$. $\det = -16 - 4 = -20 < 0 \implies$ \textcolor{eSaddleBrown}{Indefinite}.

\geminiwarning{The handwritten notes contained a matrix that was not symmetric ($\begin{smallmatrix} -5 & 0 & 0 \\ 0 & -4 & -7 \\ 3 & -2 & -4 \end{smallmatrix}$). The Hurwitz criterion strictly applies \emph{only} to symmetric matrices! Let's look at the correct symmetric matrix the TA meant to evaluate:}

\textbf{3.} $A = \begin{pmatrix} -5 & 0 & 0 \\ 0 & -4 & -2 \\ 0 & -2 & -4 \end{pmatrix}$. Let's check the Hurwitz minors:
\begin{align*}
    A_1 &= -5 \quad \textcolor{eGray}{(< 0)} \\
    A_2 &= \det\begin{pmatrix} -5 & 0 \\ 0 & -4 \end{pmatrix} = 20 \quad \textcolor{eGray}{(> 0)} \\
    A_3 &= \det(A) = -5 \cdot (16 - 4) = -60 \quad \textcolor{eGray}{(< 0)}
\end{align*}
The signs alternate ($-, +, -$), so this is \textcolor{eDarkRed}{Negative definite}.

\textbf{4.} $A = \begin{pmatrix} -1 & 0 & -4 \\ 0 & 2 & -2 \\ -4 & -2 & 22 \end{pmatrix}$. 
$A_1 = -1$ and $A_2 = \det\begin{pmatrix} -1 & 0 \\ 0 & 2 \end{pmatrix} = -2$. Since the signs are $(-, -, \dots)$, it violates both the positive and negative definite patterns $\implies$ \textcolor{eSaddleBrown}{Indefinite}.
\end{example*}

\hrulefill

\begin{exercise*}
Find the critical points of $f_\alpha(x,y) = x^3 - y^3 + 3\alpha xy$ for a constant $\alpha \in \R \setminus \set{0}$, and classify them (min/max/saddle).
\end{exercise*}

\begin{solution}
First, find the critical points by setting the gradient to zero:
\begin{align}
    \partial_x f &= 3x^2 + 3\alpha y = 0 \implies y = -\frac{x^2}{\alpha} \\
    \partial_y f &= -3y^2 + 3\alpha x = 0
\end{align}
Substitute (1) into (2):
\[ -3\left(-\frac{x^2}{\alpha}\right)^2 + 3\alpha x = 0 \implies -3\frac{x^4}{\alpha^2} + 3\alpha x = 0 \]
Divide by 3 and factor out $x$:
\[ x \left( \alpha - \frac{x^3}{\alpha^2} \right) = 0 \]
This yields two cases:
\begin{itemize}
    \item $x = 0 \implies y = 0$. Point: \textbf{$(0,0)$}.
    \item $x^3 = \alpha^3 \implies x = \alpha \implies y = -\alpha$. Point: \textbf{$(\alpha, -\alpha)$}.
\end{itemize}

Now, compute the Hessian matrix:
\[ \mathscr{H}f(x,y) = \begin{pmatrix} 6x & 3\alpha \\ 3\alpha & -6y \end{pmatrix} \]

\textbf{Test point 1: $(0,0)$}
\[ \mathscr{H}f(0,0) = \begin{pmatrix} 0 & 3\alpha \\ 3\alpha & 0 \end{pmatrix} \]
$\det = -9\alpha^2 < 0 \implies$ \textcolor{eSaddleBrown}{\textbf{Saddle Point}}.

\textbf{Test point 2: $(\alpha, -\alpha)$}
\[ \mathscr{H}f(\alpha, -\alpha) = \begin{pmatrix} 6\alpha & 3\alpha \\ 3\alpha & 6\alpha \end{pmatrix} = 3\alpha \begin{pmatrix} 2 & 1 \\ 1 & 2 \end{pmatrix} \]
The determinant is $36\alpha^2 - 9\alpha^2 = 27\alpha^2 > 0$. The trace is $12\alpha$.
\begin{itemize}
    \item If $\alpha > 0$: Trace is positive $\implies$ Positive definite $\implies$ \textcolor{eDarkGreen}{\textbf{Local Minimum}}.
    \item If $\alpha < 0$: Trace is negative $\implies$ Negative definite $\implies$ \textcolor{eDarkRed}{\textbf{Local Maximum}}.
\end{itemize}
\end{solution}